\documentclass[11pt,a4paper]{article}
\usepackage[T1]{fontenc}
\usepackage[utf8]{inputenc}
\usepackage{amsmath,amssymb,xcolor}
\usepackage[margin=25mm]{geometry}
\usepackage[colorlinks=true,urlcolor=blue]{hyperref}
\definecolor{nexorablue}{HTML}{2459B5}
\setlength{\parindent}{0pt}
\setlength{\parskip}{7pt}
\setlength{\emergencystretch}{2em}
\renewcommand{\familydefault}{\sfdefault}
\begin{document}

{\small\bfseries\color{nexorablue} NEXORA CPGE}\par
{\small MP / NEX-SLC-04 / Corrigé}\par
{\LARGE\bfseries Identifier les deux entrées d’un moteur}\par
\vspace{0.25em}\hrule\vspace{0.6em}

\section*{Élimination du courant}
En substituant et en passant dans le domaine de Laplace avec état initial nul, on obtient
\[ [R(Jp+f)+k_tk_e]\Omega=k_t U-RC_r. \]

\section*{Deux transferts}
Pour l’entrée de commande, annuler $C_r$ : le transfert est $k_t/[R(Jp+f)+k_tk_e]$. Pour la perturbation, annuler $U$ : il reste $-R/[R(Jp+f)+k_tk_e]$. Un couple résistant positif s’oppose au mouvement, d’où le signe négatif.


\end{document}
